\chapter{Series de potencias}\label{ch:b2:powerseries}

Las series de potencias son las series de funciones que mejor se
portan de toda la matemática: dentro de su disco de convergencia
convergen \omterm{def:b2:funcseq:series}{normalmente} sobre los \omterm{def:b2:metric:compact}{compactos}, pueden derivarse e
integrarse término a término sin pensárselo dos veces, y sus sumas
---las \emph{\omterm{def:b2:powerseries:analytic}{funciones analíticas}}--- quedan determinadas por sus
coeficientes. Este capítulo demuestra ese paquete completo y recupera
honestamente todos los desarrollos de Taylor del volumen del primer
año; las \omterm{ex:b2:powerseries:fibonacci}{funciones generatrices} lo cierran con dividendos
algebraicos.

\section{Radio de convergencia}

\begin{lemma}[Abel]\label{lem:b2:powerseries:abel}
Si la sucesión $(a_n z_0^n)$ está acotada para algún $z_0 \neq 0$,
entonces $\sum a_n z^n$ converge \omterm{def:b2:series:def}{absolutamente} para todo $\abs z <
\abs {z_0}$, y \omterm{def:b2:funcseq:series}{normalmente} sobre cada disco $\abs z \leq r <
\abs{z_0}$.
\end{lemma}

\begin{proof}
Con $\abs{a_n z_0^n} \leq M$ y $\abs z \leq r$:
\[
\abs{a_n z^n} = \abs{a_n z_0^n}\,\Bigl|\frac{z}{z_0}\Bigr|^n
\leq M\Bigl(\frac{r}{\abs{z_0}}\Bigr)^{\!n},
\]
una cota geométrica convergente y uniforme sobre el disco.
\end{proof}

\begin{definition}[Radio de convergencia]\label{def:b2:powerseries:radius}
El \emph{radio de convergencia}\index{radio de convergencia} de
$\sum a_n z^n$ es
\[
R = \sup\{r \geq 0 : (a_n r^n) \text{ está acotada}\} \in
\intcc{0}{+\infty} .
\]
Por el~\cref{lem:b2:powerseries:abel}: convergencia absoluta para
$\abs z < R$ (normal sobre subdiscos \omterm{def:b2:metric:compact}{compactos}) y divergencia ---de
hecho, términos no acotados--- para $\abs z > R$. Sobre la
circunferencia frontera puede ocurrir cualquier cosa
(\cref{exo:b2:powerseries:2}). En la práctica, $R$ se calcula con el
criterio del cociente de d'Alembert sobre $\abs{a_n}\abs z^n$ o por
comparación.
\end{definition}

\begin{example}[Un radio sin criterio del cociente]\label{ex:b2:powerseries:sinn}
¿Cuál es el radio de $\sum \sin(n)\,z^n$? El cociente
$\abs{\sin(n+1)/\sin n}$ no tiene límite, pero la definición
funciona directamente. \emph{$R \geq 1$:} $\abs{\sin n} \leq 1$,
luego $(\sin n\cdot r^n)$ está acotada para todo $r < 1$ ---de
hecho, para $r = 1$---. \emph{$R \leq 1$:} basta con que $\sin n
\not\to 0$. Supongamos $\sin n \to 0$; la fórmula de adición
\[
\sin(n+1) = \sin n\cos 1 + \cos n\sin 1
\]
forzaría $\cos n \to 0$ (despejando $\cos n$, pues $\sin 1 \neq
0$), en contra de $\sin^2 n + \cos^2 n = 1$. Así pues, los términos
$\sin(n)\,1^n$ no tienden a $0$: la serie diverge en $z = 1$ y $R
\leq 1$. Conclusión: $R = 1$. Moraleja: el radio es un enunciado
sobre la acotación de $\abs{a_n}r^n$; nunca hace falta ningún
límite de cocientes, y los argumentos de acotación zanjan casos que
el criterio del cociente ni siquiera toca (compárese con los
coeficientes oscilantes del~\cref{exo:b2:powerseries:1}).
\end{example}

\begin{proposition}[Operaciones]\label{prop:b2:powerseries:operations}
Sean $\sum a_nz^n$ y $\sum b_nz^n$ de radios $R_a, R_b$. Entonces,
para $\abs z < \min(R_a, R_b)$,
\[
\sum (a_n + b_n)z^n = \sum a_nz^n + \sum b_nz^n,
\qquad
\Bigl(\sum a_nz^n\Bigr)\Bigl(\sum b_nz^n\Bigr) = \sum c_n z^n,
\quad c_n = \sum_{k=0}^{n} a_kb_{n-k},
\]
y ambas series tienen radio $\geq \min(R_a, R_b)$. (El producto es
el \omterm{thm:b2:series:fubini}{producto de Cauchy}, legítimo por la convergencia absoluta y el~\cref{thm:b2:series:fubini}.)
\end{proposition}

\begin{proof}
La fórmula de la suma es la linealidad de las series convergentes, y
$(a_n + b_n)r^n$ está acotada siempre que lo estén $a_nr^n$ y
$b_nr^n$: radio $\geq \min(R_a, R_b)$. Para el producto, fijemos
$\abs z < \min(R_a, R_b)$: ambas series convergen allí
\emph{\omterm{def:b2:series:def}{absolutamente}} (\cref{lem:b2:powerseries:abel}), de modo que la
familia doblemente indexada $(a_kz^k\,b_lz^l)_{k,l}$ es \omterm{def:b2:series:summable}{sumable} y
el~\cref{thm:b2:series:fubini} autoriza cualquier agrupación.
Agrupando por $k + l = n$:
\[
\Bigl(\sum_k a_kz^k\Bigr)\Bigl(\sum_l b_lz^l\Bigr)
= \sum_{n\geq0}\Bigl(\sum_{k+l=n}a_kb_l\Bigr)z^n
= \sum_{n\geq0}c_nz^n ,
\]
\omterm{def:b2:series:def}{absolutamente} convergente para todo $z$ así: la serie producto tiene
también radio $\geq \min(R_a, R_b)$.
\end{proof}

\begin{example}[Un cuadrado de Cauchy, contrastado]\label{ex:b2:powerseries:cauchysquare}
Elevemos al cuadrado la serie geométrica: para $\abs x < 1$, el
coeficiente de $x^n$ en $\bigl(\sum x^k\bigr)^2$ es $c_n =
\sum_{k+l=n} 1\cdot1 = n + 1$, luego
\[
\frac{1}{(1-x)^2} = \sum_{n\geq0}(n+1)\,x^n .
\]
Contraste por derivación término a término
(\cref{thm:b2:powerseries:calculus} de más abajo): derivando
$\frac{1}{1-x} = \sum x^n$ resulta $\frac{1}{(1-x)^2} = \sum
nx^{n-1} = \sum(n+1)x^n$; la misma serie por dos mecanismos sin
relación. Moraleja: cuando una identidad entre coeficientes parece
misteriosa, uno de estos dos motores (convolución o derivación)
suele producirla en una línea; la pregunta del problema de fin de
semana sobre $\sum\binom{2k}k\binom{2n-2k}{n-k} = 4^n$ hace
funcionar el motor de la convolución a plena potencia.
\end{example}

\begin{example}[Multiplicar por $\frac{1}{1-x}$ suma los coeficientes]\label{ex:b2:powerseries:partialsums}
Un \omterm{thm:b2:series:fubini}{producto de Cauchy} contra la serie geométrica tiene un
significado memorable: para toda $\sum a_nx^n$ de radio $R > 0$ y
$\abs x < \min(R, 1)$,
\[
\frac{1}{1-x}\sum_{n\geq0}a_nx^n
= \sum_{n\geq0}\Bigl(\sum_{k=0}^{n}a_k\Bigr)x^n :
\]
multiplicar por $\frac{1}{1-x}$ sustituye los coeficientes por sus
sumas parciales (convolución con la sucesión de unos). Ejemplo:
$\dfrac{\eu^x}{1-x} = \sum_n s_n x^n$ con $s_n = \sum_{k\leq
n}\frac{1}{k!}$, las sumas parciales de $\eu$; compárese con el~\cref{exo:b2:powerseries:11}, donde el mismo producto con $\eu^{-x}$
codifica los recuentos de desarreglos. Moraleja: las operaciones
sobre series de potencias son operaciones sobre sucesiones de
coeficientes disfrazadas (multiplicar por $\frac1{1-x}$: sumar;
multiplicar por $x$: desplazar; derivar: multiplicar por $n$ y
desplazar), un diccionario que el capítulo de \omterm{ex:b2:powerseries:fibonacci}{funciones generatrices}
leerá con soltura.
\end{example}

\section{Regularidad de la suma}

\begin{theorem}[Cálculo término a término]\label{thm:b2:powerseries:calculus}
Sea $f(x) = \sum_{n\geq0} a_n x^n$ de radio $R > 0$ (variable real
$x \in \intoo{-R}{R}$).
\begin{enumerate}
  \item La serie derivada $\sum n\,a_n x^{n-1}$ tiene el
        \emph{mismo} radio $R$, y $f$ es de clase $C^1$ con $f'(x) =
        \sum_{n \geq 1} n a_n x^{n-1}$. Iterando, $f$ es $C^\infty$
        y
        \[
        a_n = \frac{f^{(n)}(0)}{n!} :
        \]
        los coeficientes de una serie de potencias son únicos (dos
        series con la misma suma cerca de $0$ tienen los mismos
        coeficientes).
  \item Primitiva término a término: $\sum \frac{a_n}{n+1}x^{n+1}$
        tiene radio $R$ y derivada $f$.
\end{enumerate}
\end{theorem}

\begin{proof}
\emph{Mismo radio:} si $(a_nr^n)$ está acotada y $r' < r$, entonces
$n\abs{a_n} r'^{\,n-1} = \frac{n}{r'}\abs{a_nr^n}
\bigl(\frac{r'}{r}\bigr)^n$ está acotada (de hecho, $\to 0$: lo
geométrico gana a $n$), luego $R' \geq R$; recíprocamente, $\abs{a_n
x^n} \leq \abs x \cdot n\abs{a_n}\abs x^{n-1}$ da $R \geq R'$.

\emph{Derivación:} sobre $\intcc{-r}{r}$ con $r < R$, la serie
derivada converge \omterm{def:b2:funcseq:series}{normalmente} ($n\abs{a_n}r^{n-1}$ es \omterm{def:b2:series:summable}{sumable} por
el cálculo del radio); la original converge en $x = 0$: el teorema
de derivación para series
(\cref{thm:b2:funcseq:seriestransfer}) se aplica sobre cada segmento
así y, por tanto, sobre $\intoo{-R}{R}$. Iterando $k$ veces y
evaluando en $0$: explícitamente, la serie derivada $k$-ésima es
\[
f^{(k)}(x) = \sum_{n\geq k} n(n-1)\cdots(n-k+1)\,a_n\,x^{n-k},
\]
y en $x = 0$ se anula todo término con $n > k$, quedando solo el
término constante $k(k-1)\cdots1\cdot a_k$: $f^{(k)}(0) =
k!\,a_k$. De ahí se sigue la unicidad de los coeficientes: dos
series de potencias con la misma suma cerca de $0$ tienen las mismas
derivadas en $0$ y, por tanto, los mismos $a_k$. Primitivas: mismo
radio por el mismo cálculo, y se deriva término a término de
vuelta.
\end{proof}

\begin{example}[Evaluar una serie en un punto]\label{ex:b2:powerseries:evaluate}
¿Cuánto vale $\sum_{n\geq1}\dfrac{n^2}{2^n}$? Es la suma $\sum
n^2x^n$ del~\cref{exo:b2:powerseries:3} evaluada \emph{dentro} del
disco, en $x = \frac12 < 1 = R$, donde toda manipulación empleada
para deducir la forma cerrada era legítima:
\[
\sum_{n\geq1} n^2x^n = \frac{x(1+x)}{(1-x)^3}
\quad\Longrightarrow\quad
\sum_{n\geq1}\frac{n^2}{2^n}
= \frac{\frac12\cdot\frac32}{(\frac12)^3}
= \frac{3/4}{1/8} = 6 .
\]
Mismo motor, otros mandos: $x = \frac13$ da $\sum\frac{n^2}{3^n} =
\frac{\frac13\cdot\frac43}{(2/3)^3} = \frac32$. Moraleja: una
identidad entre series de potencias es una máquina y no una sola
fórmula; una única deducción tarifa de golpe todas las series
numéricas $\sum n^2q^n$ para todo $\abs q < 1$, y así es como el
capítulo de \omterm{ex:b2:powerseries:fibonacci}{funciones generatrices} calculará esperanzas y varianzas
al por mayor.
\end{example}

\begin{example}[Los clásicos, esta vez honestamente]\label{ex:b2:powerseries:classics}
$\displaystyle\frac{1}{1 - x} = \sum x^n$ ($R = 1$); integrando
término a término (\cref{thm:b2:powerseries:calculus}~(2)):
\[
-\ln(1 - x) = \sum_{n\geq1} \frac{x^n}{n},
\qquad
\arctan x = \sum_{n \geq 0} \frac{(-1)^n x^{2n+1}}{2n+1}
\quad (\abs x < 1),
\]
la segunda en dos pasos: sustitúyase $-x^2$ en la serie geométrica
para obtener $\frac{1}{1+x^2} = \sum(-1)^nx^{2n}$ (de radio $1$, ya
que $\abs{x^2} < 1 \iff \abs x < 1$) y tómese después la primitiva
término a término que se anula en $0$; ambos miembros son primitivas
de la misma función con el mismo valor en $0$ y, por tanto, son
iguales sobre $\intoo{-1}{1}$. Y $\exp$: la serie $E(x) = \sum
\frac{x^n}{n!}$ ($R = \infty$) cumple $E' = E$ y $E(0) = 1$ por
derivación término a término, luego $E = \exp$ por la unicidad del
primer año. Todo ``desarrollo estándar'' del volumen del primer año
es ya un teorema sobre su serie de potencias completa.
\end{example}

\begin{example}[Un logaritmo calculado desde dentro del disco]\label{ex:b2:powerseries:lntwo}
Evaluando $-\ln(1-x) = \sum\frac{x^n}{n}$ en el punto interior $x =
\frac12$:
\[
\sum_{n\geq1}\frac{1}{n\,2^n} = \ln 2 ,
\]
una representación de $\ln 2$ de convergencia rápida (diez términos
dan ya $0.69306\ldots$ frente a $\ln 2 = 0.69314\ldots$), mucho
mejor que la serie alternada $1 - \frac12 + \frac13 - \dots$,
disponible solo en la frontera. Moraleja: siempre que una constante
sea alcanzable tanto en el borde como estrictamente dentro del
disco, gana numéricamente el interior: decrecimiento geométrico
contra decrecimiento armónico.
\end{example}

\begin{example}[La derivación conserva el radio, no la frontera]\label{ex:b2:powerseries:boundaryloss}
La serie $\sum_{n\geq1}\frac{x^n}{n^2}$ tiene radio $1$ y converge
en \emph{ambos} extremos ($\sum\frac1{n^2}$ y su gemela alternada).
Su serie derivada,
\[
\sum_{n\geq1}\frac{x^{n-1}}{n} ,
\]
tiene el mismo radio $1$ ---como garantiza el~\cref{thm:b2:powerseries:calculus}---, pero ahora diverge en $x = 1$
(serie armónica) mientras que sigue convergiendo en $x = -1$
(alternada). Una derivación más da $\sum_{n\geq2}\frac{n-1}{n}
x^{n-2}$, divergente en ambos extremos (los términos no tienden a
$0$). Moraleja: cada derivación multiplica los coeficientes por $n$,
lo que nunca mueve el radio (lo geométrico gana a lo polinómico)
pero se come un orden de decrecimiento en la frontera; el cálculo
término a término es un deporte de interior, y lo que ocurra en el
borde ha de reexaminarse: la teoría de Abel--Tauber del problema de
fin de semana es exactamente ese reexamen.
\end{example}

\begin{example}[Separar una serie por restos ---hasta el final]\label{ex:b2:powerseries:quarter}
Calculemos $f(x) = \sum_{n\geq0} \dfrac{x^{4n}}{(4n)!}$ en forma
cerrada. Tanto $\cosh x = \sum \frac{x^{2m}}{(2m)!}$ como $\cos x =
\sum \frac{(-1)^m x^{2m}}{(2m)!}$ tienen radio $\infty$, de modo que
su media puede calcularse término a término:
\[
\frac{\cosh x + \cos x}{2}
= \sum_{m\geq0}\frac{1 + (-1)^m}{2}\,\frac{x^{2m}}{(2m)!}
= \sum_{m \text{ par}}\frac{x^{2m}}{(2m)!}
= \sum_{n\geq0}\frac{x^{4n}}{(4n)!} = f(x) .
\]
El filtro $\frac{1+(-1)^m}{2}$ retiene exactamente los $m$ pares:
este es el avatar real del filtro por raíces de la unidad (la
versión compleja, con $\iu^n$, extrae los restos módulo $4$ de un
golpe). Comprobación final: $f$ resuelve $f'''' = f$ con $f(0) = 1$
y $f'(0) = f''(0) = f'''(0) = 0$ ---derívese la serie cuatro veces
(\cref{thm:b2:powerseries:calculus}) y véase cómo se reproduce---; y
$\frac{\cosh + \cos}{2}$ satisface los mismos datos.
\end{example}

\begin{definition}[Funciones analíticas]\label{def:b2:powerseries:analytic}
$f$ es \emph{analítica}\index{función analítica} en $x_0$ cuando es
la suma de una serie de potencias en $(x - x_0)$ sobre un entorno; y
sobre un intervalo, cuando lo es en todo punto. Las sumas de series
de potencias son analíticas dentro de su disco (por reordenación del
desarrollo, admitida a este nivel para el recentrado; el caso $x_0 =
0$ es el~\cref{thm:b2:powerseries:calculus}). Analítica implica
$C^\infty$; el recíproco \emph{falla}: la función plana
$\eu^{-1/x^2}$ (\cref{exo:b2:powerseries:7}).
\end{definition}

\begin{example}[Recentrar, y el radio como distancia]\label{ex:b2:powerseries:recenter}
Desarrollemos $f(x) = \frac{1}{1-x}$ alrededor de $x_0 = \frac12$:
escribiendo $x = \frac12 + h$,
\[
\frac{1}{1 - x} = \frac{1}{\frac12 - h}
= \frac{2}{1 - 2h}
= \sum_{n\geq0} 2^{n+1}\,h^n
= \sum_{n\geq0} 2^{n+1}\Bigl(x - \frac12\Bigr)^{\!n},
\]
válido para $\abs{2h} < 1$, es decir, $\abs{x - \frac12} <
\frac12$. El nuevo radio es exactamente la distancia del nuevo
centro a la singularidad $x = 1$: recentrar encoge (o agranda) el
disco hasta ajustarlo a la obstrucción más próxima. Moraleja: esta
es la imagen que hay tras la definición de analiticidad ---una
función, muchas series de potencias locales, cada una viviendo en el
mayor disco que evita el problema---; el volumen del tercer año
convierte la heurística ``radio $=$ distancia a la singularidad
compleja más próxima'' en un teorema.
\end{example}

\begin{remark}[Errores frecuentes]
\emph{(i) El criterio del cociente es suficiente, no necesario:}
cuando $\abs{a_{n+1}/a_n}$ no tiene límite
(\cref{ex:b2:powerseries:sinn}, \cref{exo:b2:powerseries:1}), hay
que volver a la definición: $R = \sup\{r : (a_nr^n)$ acotada$\}$.
\emph{(ii) Nada cruza la frontera gratis:} la derivación y la
integración término a término son teoremas \emph{dentro} del disco
\omterm{def:b2:metric:topology}{abierto}; en $\abs x = R$, cada serie ha de reexaminarse (ese es todo
el tema del problema de fin de semana). \emph{(iii) Radio de una
suma:} $\min(R_a, R_b)$ es solo una cota inferior, pues las
cancelaciones pueden agrandarlo ($a_n = 1$, $b_n = -1$: suma
idénticamente $0$, radio $\infty$). \emph{(iv) $C^\infty$ no es
\omterm{def:b2:powerseries:analytic}{analítica}:} una serie de Taylor convergente puede converger a la
función \emph{equivocada} (\cref{exo:b2:powerseries:7}); antes de
escribir $f(x) = \sum \frac{f^{(n)}(0)}{n!}x^n$ hay que
demostrarlo, vía una ecuación diferencial
(\cref{met:b2:powerseries:ode}), una estimación del resto o una
fórmula integral.
\end{remark}

\begin{remark}[Dónde se usa]
Las series de potencias son el caballo de batalla de tres capítulos
posteriores: el de ecuaciones diferenciales resuelve ecuaciones
lineales inyectando $\sum a_nx^n$ (el recuadro de método de más
abajo, industrializado); el de \omterm{ex:b2:powerseries:fibonacci}{funciones generatrices} convierte
identidades sobre probabilidades en identidades sobre series de
potencias y viceversa; y el volumen del tercer año oficializa la
variable compleja, donde la analiticidad pasa a ser equivalente a la
derivabilidad compleja y el ``recentrado admitido'' de más arriba
recibe su demostración honesta. El problema de fin de semana explora
el único lugar donde los teoremas de este capítulo callan: la propia
frontera $\abs x = R$.
\end{remark}

\begin{method}[Desarrollar mediante una ecuación diferencial]\label{met:b2:powerseries:ode}
Para desarrollar una función $f$ en serie de potencias: hállese una
ecuación diferencial lineal con coeficientes polinómicos que
satisfaga $f$; inyéctese $\sum a_nx^n$; identifíquense coeficientes
para obtener una recurrencia sobre $(a_n)$; resuélvase y
compruébense el radio y las condiciones iniciales. Ejemplo ---la
serie binomial---: $f(x) = (1+x)^\alpha$ satisface $(1+x)f' = \alpha
f$ con $f(0) = 1$; al inyectar resulta $(n+1)a_{n+1} = (\alpha -
n)a_n$, luego $a_n = \binom{\alpha}{n}$, con radio $1$ (criterio del
cociente), y la suma, que satisface la misma ecuación con el mismo
valor inicial, vale $(1 + x)^\alpha$ por el teorema de unicidad para
ecuaciones diferenciales lineales (volumen del primer año).
\end{method}

\begin{example}[El método sobre una ecuación con término forzante]\label{ex:b2:powerseries:odeforced}
Resolvamos $y' = y + x$, $y(0) = 0$, por series de potencias.
Inyectando $y = \sum a_nx^n$:
\[
\sum_{n\geq0}(n+1)a_{n+1}x^n
= \sum_{n\geq0}a_nx^n + x ,
\]
e identificando coeficientes: $a_1 = a_0 = 0$, $2a_2 = a_1 + 1 = 1$
y $(n+1)a_{n+1} = a_n$ para $n \geq 2$. Así pues, $a_2 =
\frac{1}{2!}$ y, por inducción, $a_n = \frac{1}{n!}$ para todo $n
\geq 2$: radio $\infty$, y
\[
y(x) = \sum_{n\geq2}\frac{x^n}{n!} = \eu^x - 1 - x .
\]
Comprobación: $y' = \eu^x - 1 = y + x$ y $y(0) = 0$. Moraleja: la
recurrencia \emph{es} la ecuación, coeficiente a coeficiente; el
término forzante solo perturba un número finito de coeficientes
iniciales, tras los cuales toma el relevo el patrón homogéneo: una
sombra discreta de ``solución particular más solución homogénea''.
\end{example}

\section{Funciones generatrices}

\begin{example}[Fibonacci]\label{ex:b2:powerseries:fibonacci}
Sea $F(x) = \sum_{n\geq0} F_n x^n$ (números de Fibonacci, con $F_0 =
0$ y $F_1 = 1$). La recurrencia $F_{n+2} = F_{n+1} + F_n$ se
traduce, multiplicando por $x^{n+2}$ y sumando, en
\[
F(x) - x = x\,F(x) + x^2 F(x)
\quad\Longrightarrow\quad
F(x) = \frac{x}{1 - x - x^2} ,
\]
válido donde la serie converge. El radio es $\frac{1}{\varphi}$: de
$F_n \sim \frac{\varphi^n}{\sqrt5}$ (Binet, ejemplo siguiente; o
bien la inducción tosca $F_n \leq 2^n$ más la recurrencia), el
criterio del cociente da
\[
\frac{F_{n+1}\abs x^{n+1}}{F_n\abs x^n}
\longrightarrow \varphi\abs x ,
\qquad\text{convergencia si y solo si } \abs x < \frac1\varphi
\approx 0.618 .
\] La descomposición en fracciones simples de
$\frac{x}{1 - x - x^2}$ y la serie geométrica vuelven a deducir la
fórmula de Binet: las \omterm{ex:b2:powerseries:fibonacci}{funciones generatrices} industrializan las
recurrencias lineales.\index{función generatriz}
\end{example}

\begin{example}[La fórmula de Binet, ejecutada]\label{ex:b2:powerseries:binet}
Sean $\varphi = \frac{1+\sqrt5}{2}$ y $\psi = \frac{1-\sqrt5}{2}$,
las raíces de $X^2 = X + 1$; como $\varphi + \psi = 1$ y
$\varphi\psi = -1$,
\[
1 - x - x^2 = (1 - \varphi x)(1 - \psi x) .
\]
Fracciones simples: buscando $\frac{x}{(1-\varphi x)(1-\psi x)} =
\frac{A}{1 - \varphi x} + \frac{B}{1 - \psi x}$, el término
constante da $A + B = 0$ y el coeficiente en $x$ da $-A\psi -
B\varphi = 1$, luego $A(\varphi - \psi) = 1$: $A =
\frac{1}{\sqrt5} = -B$. Dos series geométricas después,
\[
F(x) = \frac{1}{\sqrt5}\sum_{n\geq0}
\bigl(\varphi^n - \psi^n\bigr)x^n
\quad\Longrightarrow\quad
F_n = \frac{\varphi^n - \psi^n}{\sqrt5}
\]
por la unicidad de los coeficientes
(\cref{thm:b2:powerseries:calculus}). Como $\abs\psi < 1$, el
término $\frac{\psi^n}{\sqrt5}$ tiene valor absoluto $< \frac12$:
$F_n$ es el \emph{entero más próximo} a $\frac{\varphi^n}{\sqrt5}$.
Moraleja: el radio $\frac1\varphi$ de $F$ es el inverso de la raíz
dominante; el crecimiento de los coeficientes y el
\omterm{def:b2:powerseries:radius}{radio de convergencia} son la misma información leída en
direcciones opuestas.
\end{example}

\begin{example}[Números de Catalan]\label{ex:b2:powerseries:catalan}
Los \omterm{ex:b2:powerseries:catalan}{números de Catalan} $C_n$ (número de triangulaciones, de
formas de poner paréntesis, de caminos de Dyck, \dots) cumplen $C_0
= 1$ y $C_{n+1} = \sum_{k=0}^n C_kC_{n-k}$. La
\omterm{ex:b2:powerseries:fibonacci}{función generatriz} $C(x) = \sum C_nx^n$ satisface entonces
(¡\omterm{thm:b2:series:fubini}{producto de Cauchy}!)
\[
C(x) = 1 + x\,C(x)^2
\quad\Longrightarrow\quad
C(x) = \frac{1 - \sqrt{1 - 4x}}{2x} ,
\]
eligiendo la raíz con $C(0) = 1$: resolver la cuadrática $xC^2 - C +
1 = 0$ da los dos candidatos $\frac{1 \pm \sqrt{1-4x}}{2x}$ y,
cuando $x \to 0$, la raíz con ``$+$'' explota como $\frac1x$
mientras que la de ``$-$'' tiende a $1$ (desarróllese $\sqrt{1-4x} =
1 - 2x + O(x^2)$): solo el signo menos puede llevar una serie de
potencias con $C_0 = 1$. Desarrollando $\sqrt{1 - 4x}$ con la serie
binomial se obtiene la forma cerrada
\[
C_n = \frac{1}{n+1}\binom{2n}{n} ,
\]
ejecutada en el~\cref{exo:b2:powerseries:8}.\index{números de
Catalan}
\end{example}

\begin{remark}[Series formales frente a series convergentes]
Todo cálculo con \omterm{ex:b2:powerseries:fibonacci}{funciones generatrices} de más arriba termina
invocando la unicidad de los coeficientes, y ese teorema vive
\emph{dentro} de un disco de radio positivo: antes de ``leer'' $F_n$
o $C_n$ hay que saber que $R > 0$. Basta una cota tosca a priori:
$F_n \leq 2^n$ (inducción inmediata) da $R \geq \frac12$ para
Fibonacci; $C_n \leq 4^n$ (cada \omterm{ex:b2:powerseries:catalan}{número de Catalan} cuenta
subconjuntos de caminos) da $R \geq \frac14$. Cuidado con el extremo
degenerado de la escala: $\sum n!\,x^n$ tiene radio $0$, y
manipularla como una función carece de sentido; las identidades que
involucran series así pertenecen al cálculo \emph{formal} de
coeficientes, un juego puramente algebraico con sus propias reglas
(distintas). A este nivel: asegúrese siempre primero un radio
positivo y calcúlese después con libertad dentro de él.
\end{remark}

\begin{remark}[Perspectivas dentro de este volumen]
Las series de potencias son una de las dos grandes máquinas de
desarrollo del libro; la otra son las series de Fourier de los
capítulos armónicos, y compararlas resulta instructivo. Una serie de
potencias es rígida: sus coeficientes están forzados ($a_n =
f^{(n)}(0)/n!$), su convergencia es implacable (normal dentro,
imposible fuera) y su suma es \omterm{def:b2:powerseries:analytic}{analítica}, es decir, infinitamente
rígida (\cref{def:b2:powerseries:analytic}). Una serie de Fourier es
flexible: representa incluso señales meramente suaves a trozos, al
precio de delicadas cuestiones de convergencia en la frontera de la
suavidad. Las dos teorías se encuentran en el problema de fin de
semana de este capítulo: la sumación de Cesàro y la \omterm{thm:b2:series:abel}{sumación de
Abel}, desarrolladas aquí para la circunferencia frontera, regresan
en el capítulo de Fourier como los núcleos de Fejér y de Poisson.
Entretanto, el capítulo de ecuaciones diferenciales consume series
de potencias directamente ($\eu^{tA}$, soluciones en serie), y el de
\omterm{ex:b2:powerseries:fibonacci}{funciones generatrices} convierte el truco del~\cref{ex:b2:powerseries:fibonacci} en un cálculo sistemático para
probabilidades.
\end{remark}

\section{Ejercicios}

\begin{exercise}[$\star$]\label{exo:b2:powerseries:1}
\omterm{def:b2:powerseries:radius}{Radios de convergencia}: $\sum \dfrac{n^2}{2^n}z^n$; $\;\sum
\dfrac{z^n}{\binom{2n}{n}}$; $\;\sum z^{n!}$; $\;\sum \bigl(2 +
(-1)^n\bigr)^n z^n$.
\end{exercise}

\begin{exercise}[$\star$]\label{exo:b2:powerseries:2}
Prueba que $\sum z^n$, $\sum \frac{z^n}{n}$ y $\sum
\frac{z^n}{n^2}$ tienen todas radio $1$ pero se comportan de manera
distinta en $z = 1$ y $z = -1$: divergencia/divergencia,
divergencia/convergencia, convergencia/convergencia.
\end{exercise}

\begin{exercise}[$\star$]\label{exo:b2:powerseries:3}
Calcula las sumas, para $\abs x < 1$:
\[
\sum_{n\geq0} n x^n,
\qquad
\sum_{n\geq0} n^2 x^n,
\qquad
\sum_{n\geq0} \frac{x^{2n+1}}{2n+1} .
\]
\end{exercise}

\begin{exercise}[$\star\star$]\label{exo:b2:powerseries:4}
Desarrolla en serie de potencias en $0$, con su radio:
$\dfrac{1}{(1-x)(2-x)}$ (fracciones simples); $\;\ln(1 + x + x^2)$
\emph{(escribe $1 + x + x^2 = \frac{1 - x^3}{1 - x}$)}.
\end{exercise}

\begin{exercise}[$\star\star$]\label{exo:b2:powerseries:5}
Demuestra que $f(x) = \sum_{n\geq1} H_n x^n = -\dfrac{\ln(1 - x)}{1
- x}$ para $\abs x < 1$, donde $H_n$ es el número armónico
\emph{(\omterm{thm:b2:series:fubini}{producto de Cauchy} de $\sum x^n$ y $\sum \frac{x^n}{n}$)}.
\end{exercise}

\begin{exercise}[$\star\star$]\label{exo:b2:powerseries:6}
Resuelve mediante \omterm{ex:b2:powerseries:fibonacci}{función generatriz} la recurrencia $u_0 = 1$,
$u_{n+1} = 2u_n + n$: calcula $U(x) = \sum u_nx^n$ en forma cerrada,
descompón y lee $u_n = 2^{n+1} - n - 1$.
\end{exercise}

\begin{exercise}[$\star\star$]\label{exo:b2:powerseries:7}
Sea $f(x) = \eu^{-1/x^2}$ para $x \neq 0$ y $f(0) = 0$. Demuestra
que $f$ es de clase $C^\infty$ sobre $\R$ con $f^{(n)}(0) = 0$ para
todo $n$ \emph{(prueba por inducción que $f^{(n)}(x) =
P_n\bigl(\frac1x\bigr) \eu^{-1/x^2}$ para ciertos polinomios $P_n$,
y usa la comparación de crecimientos)}. Concluye que $f$ no es
\omterm{def:b2:powerseries:analytic}{analítica} en $0$: su serie de Taylor en $0$ converge, pero a la
función equivocada.
\end{exercise}

\begin{exercise}[$\star\star\star$]\label{exo:b2:powerseries:8}
Completa el~\cref{ex:b2:powerseries:catalan}: desarrolla $\sqrt{1 -
4x}$ con la serie binomial, probando que
\[
\binom{1/2}{n+1}(-4)^{n+1} = -\frac{2}{n+1}\binom{2n}{n},
\]
y deduce $C_n = \frac{1}{n+1}\binom{2n}{n}$; determina el
\omterm{def:b2:powerseries:radius}{radio de convergencia} de $C(x)$ y la asintótica de $C_n$ mediante
Stirling.
\end{exercise}

\begin{exercise}[$\star\star\star$]\label{exo:b2:powerseries:9}
(Teorema del límite radial de Abel, caso particular) Supongamos que
$\sum a_n$ converge. Demuestra que $\lim_{x \to 1^-} \sum_{n} a_n
x^n = \sum_n a_n$. \emph{(\omterm{thm:b2:series:abel}{Sumación de Abel}: con $A_n$ las sumas
parciales y $A = \lim A_n$, escribe $\sum a_nx^n = (1 - x)\sum A_n
x^n$; después $\sum a_nx^n - A = (1-x)\sum (A_n - A)x^n$, y separa
la suma en un $N$ grande.)} Aplicación: $\sum \frac{(-1)^{n-1}}{n} =
\ln 2$ y $\sum \frac{(-1)^n}{2n+1} = \frac\pi4$, vueltos a demostrar
a partir de la serie de potencias.
\end{exercise}

\begin{exercise}[$\star$]\label{exo:b2:powerseries:10}
Prueba que $\displaystyle\sum_{n\geq1}\frac{x^n}{n(n+1)} = 1 +
\frac{1-x}{x}\,\ln(1-x)$ para $0 < \abs x < 1$, determina el radio y
comprueba que la convergencia es normal sobre $\intcc{-1}{1}$;
verifica que el valor en $x = 1$ predicho por la \omterm{def:b2:metric:continuity}{continuidad}
concuerda con la suma telescópica $\sum \frac{1}{n(n+1)} = 1$.
\end{exercise}

\begin{exercise}[$\star\star$]\label{exo:b2:powerseries:11}
(Desarreglos) Sea $D_n$ el número de permutaciones de $n$ objetos
sin punto fijo ($D_0 = 1$). Clasificando las permutaciones de $\{1,
\dots, n\}$ según su conjunto de puntos fijos resulta $n! =
\sum_{k=0}^{n}\binom nk D_{n-k}$. Multiplica por
$\frac{x^n}{n!}$, suma y reconoce un \omterm{thm:b2:series:fubini}{producto de Cauchy} para obtener
la \omterm{ex:b2:powerseries:fibonacci}{función generatriz} exponencial
\[
\sum_{n\geq0} D_n\,\frac{x^n}{n!} = \frac{\eu^{-x}}{1-x}
\qquad (\abs x < 1),
\]
y lee después la forma cerrada $\dfrac{D_n}{n!} =
\sum_{k=0}^{n}\dfrac{(-1)^k}{k!}$ y el límite $\dfrac{D_n}{n!} \to
\eu^{-1}$.
\end{exercise}

\begin{exercise}[$\star\star\star$]\label{exo:b2:powerseries:12}
Demuestra, con la serie binomial del~\cref{met:b2:powerseries:ode}, que
\[
\frac{1}{\sqrt{1 - 4x}} = \sum_{n\geq0}\binom{2n}{n}x^n
\qquad \Bigl(\abs x < \frac14\Bigr),
\]
y deduce, elevando al cuadrado (\omterm{thm:b2:series:fubini}{producto de Cauchy} contra
$\frac{1}{1-4x} = \sum 4^nx^n$), la identidad de convolución
\[
\sum_{k=0}^{n}\binom{2k}{k}\binom{2n-2k}{n-k} = 4^n .
\]
\end{exercise}

\section{Problema: Abel, Tauber y la frontera de la convergencia}

\begin{problem}\label{pb:b2:powerseries:1}
Dentro del disco de convergencia todo es fácil; todo el drama de las
series de potencias ocurre \emph{sobre} la frontera. Este problema
construye la teoría de frontera en la variable real: el teorema de
Abel en su forma uniforme, su recíproco bajo la condición de Tauber,
la jerarquía de métodos de sumación de Cesàro y Abel (con el teorema
de Frobenius), la integración término a término hasta la frontera
con constantes clásicas como dividendos y, por último, la rigidez de
las \omterm{def:b2:powerseries:analytic}{funciones analíticas}: el teorema de
identidad.\index{teorema de Abel}\index{teorema de Tauber} En todo
él, $(a_n)$ es una sucesión real, $f(x) = \sum_{n\geq0} a_nx^n$ y
$A_n = a_0 + \dots + a_n$.

\medskip
\noindent\textbf{Parte I --- El teorema de Abel, \omterm{def:b2:funcseq:def}{uniformemente}.}
Supongamos en esta parte que $\sum a_n$ converge, y pongamos $r_n =
\sum_{k\geq n} a_k$ (de modo que $r_n \to 0$ y $a_n = r_n -
r_{n+1}$).
\begin{enumerate}
  \item Demuestra, mediante sumación por partes, que para todos $0
        \leq x \leq 1$ y $N \leq M$,
        \[
        \Bigl|\sum_{n=N}^{M} a_n x^n\Bigr|
        \leq 2\sup_{n \geq N}\,\abs{r_n} .
        \]
  \item Deduce que $\sum a_nx^n$ converge \emph{\omterm{def:b2:funcseq:def}{uniformemente}}
        sobre $\intcc{0}{1}$, que su suma es allí \omterm{def:b2:metric:continuity}{continua}, y
        recupera el límite radial del~\cref{exo:b2:powerseries:9}:
        $f(x) \to \sum a_n$ cuando $x \to 1^-$.
  \item (Teorema de Abel para \omterm{thm:b2:series:fubini}{productos de Cauchy}) Sean $\sum a_n =
        A$ y $\sum b_n = B$, y supongamos que el \omterm{thm:b2:series:fubini}{producto de Cauchy}
        $\sum c_n$, con $c_n = \sum_{k} a_kb_{n-k}$,
        \emph{converge}, de suma $C$. Demuestra que $C = AB$
        \emph{(dentro del disco vale la identidad del producto por
        la~\cref{prop:b2:powerseries:operations}; hágase $x \to
        1^-$)}.
  \item Prueba que la hipótesis importa: para $a_n = b_n =
        \frac{(-1)^n}{\sqrt{n+1}}$, ambas series convergen y sin
        embargo $\abs{c_n} \geq \frac{2(n+1)}{n+2} \geq 1$
        \emph{(acota cada factor $\sqrt{(k+1)(n-k+1)}$ por la
        desigualdad entre medias)}: el \omterm{thm:b2:series:fubini}{producto de Cauchy} de dos
        series convergentes puede diverger.
  \item (Un dividendo del~\cref{exo:b2:powerseries:5}) Prueba que
        $\bigl(\ln(1-x)\bigr)^2 = 2\sum_{n\geq1}
        \frac{H_n}{n+1}x^{n+1}$ sobre $\intoo{-1}{1}$, comprueba
        que $\bigl(\frac{H_n}{n+1}\bigr)_{n\geq1}$ decrece a $0$ y
        concluye con Abel:
        \[
        \sum_{n\geq1} (-1)^{n+1}\,\frac{H_n}{n+1}
        = \frac{(\ln 2)^2}{2} .
        \]
\end{enumerate}

\medskip
\noindent\textbf{Parte II --- El recíproco de Tauber.} Diremos que
$\sum a_n$ es \emph{\omterm{def:b2:series:summable}{Abel-sumable}} a $L$ cuando $f(x) \to L$ al
tender $x \to 1^-$.
\begin{enumerate}[resume]
  \item Prueba que $\sum (-1)^n$ es \omterm{def:b2:series:summable}{Abel-sumable} a $\frac12$ y sin
        embargo divergente: el teorema de Abel no tiene recíproco
        incondicional.
  \item (Lema de Cesàro) Si $u_n \to 0$, entonces $\frac{u_1 +
        \dots + u_N}{N} \to 0$ \emph{(sepárese la suma en un $m$
        fijo)}.
  \item Supongamos ahora $n\,a_n \to 0$ y $f(x) \to L$. Con $x_N = 1
        - \frac1N$, demuestra las dos estimaciones
        \[
        \Bigl|\sum_{n=0}^{N} a_n\bigl(1 - x_N^n\bigr)\Bigr|
        \leq \frac{1}{N}\sum_{n=1}^{N} n\,\abs{a_n},
        \qquad
        \Bigl|\sum_{n>N} a_n x_N^n\Bigr|
        \leq \sup_{n>N}\bigl(n\abs{a_n}\bigr)
        \]
        \emph{(para la primera, $1 - x^n \leq n(1-x)$; para la
        segunda, $\abs{a_n} \leq \frac{1}{N}\sup_{m>N} m\abs{a_m}$ y
        $\sum x_N^n \leq N$)}.
  \item Concluye el \emph{teorema de Tauber}: si $n\,a_n \to 0$ y
        $\sum a_n$ es \omterm{def:b2:series:summable}{Abel-sumable} a $L$, entonces $\sum a_n$
        converge a $L$.
  \item (El tauberiano fácil para coeficientes positivos) Si $a_n
        \geq 0$ y $f$ está acotada sobre $\intco{0}{1}$, prueba que
        $\sum a_n$ converge y que $\sum a_n = \lim_{x\to1^-} f(x)$
        \emph{(acota $\sum_{n\leq N}a_nx^n \leq f(x)$, haz $x \to
        1^-$ y usa después Abel)}.
\end{enumerate}

\medskip
\noindent\textbf{Parte III --- Medias de Cesàro y teorema de
Frobenius.} Diremos que $\sum a_n$ es \emph{\omterm{def:b2:series:summable}{Cesàro-sumable}} a $L$
cuando $\sigma_N = \frac{A_0 + \dots + A_{N-1}}{N} \to L$.
\begin{enumerate}[resume]
  \item Prueba que una serie convergente es \omterm{def:b2:series:summable}{Cesàro-sumable} a su suma
        \emph{(pregunta 7 aplicada a $A_n - L$)}.
  \item Calcula el valor de Cesàro de $\sum(-1)^n$ y comprueba que
        coincide con el valor de Abel $\frac12$ de la pregunta 6.
  \item Con $S_n = A_0 + \dots + A_n = (n+1)\,\sigma_{n+1}$,
        demuestra las dos identidades, para $0 \leq x < 1$:
        \[
        f(x) = (1-x)^2\sum_{n\geq0}(n+1)\,\sigma_{n+1}x^n,
        \qquad
        (1-x)^2\sum_{n\geq0}(n+1)x^n = 1 .
        \]
  \item (Frobenius) Deduce que si $\sigma_N \to L$ entonces $f(x)
        \to L$ cuando $x \to 1^-$: \omterm{def:b2:series:summable}{Cesàro-sumable} implica
        \omterm{def:b2:series:summable}{Abel-sumable}, con el mismo valor \emph{(resta las dos
        identidades y separa la suma en un $N$ grande, como en el~\cref{exo:b2:powerseries:9})}.
  \item Prueba que la jerarquía
        \[
        \text{convergente} \;\Longrightarrow\;
        \text{Cesàro-sumable} \;\Longrightarrow\;
        \text{Abel-sumable}
        \]
        es estricta en ambas flechas: la pregunta 6 para la primera;
        para la segunda, prueba que $\sum(-1)^n(n+1)$ es
        \omterm{def:b2:series:summable}{Abel-sumable} a $\frac14$ (calcula $f$) pero no
        \omterm{def:b2:series:summable}{Cesàro-sumable} (calcula $\sigma_N$ por separado para $N$ par
        e impar).
\end{enumerate}

\medskip
\noindent\textbf{Parte IV --- Integrar hasta la frontera.}
\begin{enumerate}[resume]
  \item Supongamos que $\sum a_nx^n$ converge sobre $\intco{0}{1}$ y
        que $\sum \frac{a_n}{n+1}$ converge. Demuestra que la
        \omterm{def:b2:integration:improper}{integral impropia} $\int_0^1 f$ existe y que
        \[
        \int_0^1 \Bigl(\sum_{n\geq0} a_nx^n\Bigr)\dd x
        = \sum_{n\geq0}\frac{a_n}{n+1}
        \]
        \emph{(la primitiva $F(x) = \sum\frac{a_n}{n+1}x^{n+1}$ es
        \omterm{def:b2:metric:continuity}{continua} en $1$ por la parte I)}.
  \item Sea $\eta = \sum_{n\geq1}\frac{(-1)^{n-1}}{n^2}$. Prueba que
        $\int_0^1 \frac{\ln(1+x)}{x}\dd x = \eta$ y, separando
        índices pares e impares en la serie \omterm{def:b2:series:def}{absolutamente}
        convergente $\sum \frac1{n^2}$, que $\eta =
        \frac12\sum_{n\geq1}\frac{1}{n^2}$. (El problema de fin de
        semana del capítulo de Fourier evalúa $\sum\frac1{n^2} =
        \frac{\pi^2}{6}$.)
  \item Demuestra que
        \[
        \sum_{n\geq0}\frac{(-1)^n}{3n+1}
        = \int_0^1\frac{\dd x}{1+x^3}
        = \frac13\Bigl(\ln 2 + \frac{\pi}{\sqrt3}\Bigr)
        \]
        \emph{(la serie converge por Leibniz; intégrese la serie
        geométrica $\sum(-1)^nx^{3n}$ con la pregunta 16; después,
        fracciones simples: $\frac{1}{1+x^3} = \frac{1/3}{1+x} +
        \frac{(2-x)/3}{x^2-x+1}$)}.
  \item A partir de la serie binomial de $(1-t)^{-1/2}$
        (\cref{exo:b2:powerseries:12}), deduce
        \[
        \arcsin x = \sum_{n\geq0}
        \frac{\binom{2n}{n}}{4^n(2n+1)}\,x^{2n+1}
        \quad(\abs x < 1),
        \qquad\text{y después}\qquad
        \sum_{n\geq0}\frac{\binom{2n}{n}}{4^n(2n+1)}
        = \frac\pi2 ,
        \]
        justificando el valor en la frontera mediante la
        convergencia \emph{normal} sobre $\intcc{-1}{1}$ (úsese
        $\binom{2n}n4^{-n} \sim \frac{1}{\sqrt{\pi n}}$,
        \cref{ex:b2:comparison:centralbinomial}); aquí ni siquiera
        hace falta Abel.
  \item (Catalan en la frontera) Prueba que $\sum C_n 4^{-n} = 2$:
        la serie de Catalan del~\cref{ex:b2:powerseries:catalan}
        converge \emph{en} su radio $\frac14$ (asintótica del~\cref{exo:b2:powerseries:8}), su suma es \omterm{def:b2:metric:continuity}{continua} sobre
        $\intcc{0}{\frac14}$ y allí la forma cerrada tiene límite
        $2$.
\end{enumerate}

\medskip
\noindent\textbf{Parte V --- Rigidez: el teorema de identidad.}
\begin{enumerate}[resume]
  \item (Ceros aislados) Sea $f = \sum a_nx^n$ de radio $R > 0$ con
        no todos los $a_n$ nulos, y sea $m$ el menor índice con
        $a_m \neq 0$. Prueba que $f(x) = x^m g(x)$ con $g$ una serie
        de potencias de radio $R$ y $g(0) = a_m \neq 0$, y deduce
        que $f$ no tiene ningún cero en algún entorno perforado de
        $0$.
  \item (Teorema de identidad) Sean $f, h$ sumas de series de
        potencias cerca de $0$ y $(x_k)$ una sucesión de puntos
        \emph{no nulos} con $x_k \to 0$ y $f(x_k) = h(x_k)$.
        Demuestra que $f$ y $h$ tienen los mismos coeficientes y,
        por tanto, coinciden cerca de $0$.
  \item Halla \emph{todas} las funciones $f$ \omterm{def:b2:powerseries:analytic}{analíticas} cerca de
        $0$ con
        \[
        f\Bigl(\frac1k\Bigr) = \frac{k^2}{k^2+1}
        \qquad\text{para todo entero } k \text{ grande.}
        \]
  \item Prueba que una \omterm{def:b2:powerseries:analytic}{función analítica} sobre un intervalo
        \omterm{def:b2:metric:topology}{abierto} $I$ que se anula sobre un subintervalo se anula
        idénticamente sobre $I$ \emph{(el conjunto de los puntos en
        torno a los cuales $f$ se anula idénticamente es \omterm{def:b2:metric:topology}{abierto} y,
        por el teorema de identidad aplicado en los puntos de
        acumulación, cerrado en $I$)}. Concluye que ninguna
        \omterm{def:b2:powerseries:analytic}{función analítica} no nula sobre $\R$ tiene soporte
        \omterm{def:b2:metric:compact}{compacto}, mientras que sí existen funciones meseta de clase
        $C^\infty$ (el~\cref{exo:b2:powerseries:7} proporciona la
        pieza básica): la analiticidad es rígida y la suavidad es
        blanda.
  \item Síntesis. Una frase para cada punto: (i) qué añade el
        teorema de Abel al paquete de \omterm{def:b2:funcseq:series}{convergencia normal} del~\cref{lem:b2:powerseries:abel}; (ii) bajo qué hipótesis
        exactas vale el recíproco (Tauber) y cuál es el peldaño
        intermedio (Frobenius); (iii) una constante de frontera de
        la parte IV que ya podrías deducirle a un amigo en dos
        líneas; (iv) dónde reaparecerán las medias de Cesàro en este
        libro, para series de naturaleza muy distinta.
\end{enumerate}
\end{problem}
